% ============================================================
% Examen Final Integrador - Emprendimientos Productivos
% Transcripción literal del original (2 páginas).
% Requiere en el preámbulo: \usepackage{tikz} \usepackage{geometry}

% ============================================================
\newcommand{\exPuntos}[1]{\makebox[#1]{\dotfill}}
\newcommand{\exItem}[2]{\noindent\hspace*{18pt}\makebox[18pt][l]{#1}#2\par}
\newcommand{\exSub}[2]{\noindent\hspace*{54pt}\makebox[17pt][l]{#1}#2\par}

\clearpage
% Medidas del original (carta) escaladas a A4 (factor 0,972)
\newgeometry{left=2.12cm,right=1.17cm,top=1cm,bottom=2cm,nohead,nofoot}
\begingroup
\fontfamily{lmr}\bfseries\fontsize{10.7}{13}\selectfont
\setlength{\parindent}{0pt}
\setlength{\parskip}{0pt}

% ---------------------------------------------------------- página 1
\thispagestyle{empty}
\begin{tikzpicture}[remember picture,overlay]
  \node[anchor=north east,inner sep=0pt]
    at ([xshift=-0.9cm,yshift=-0.95cm]current page.north east)
    {\includegraphics[width=6.3cm]{logos/logo-impa}};
\end{tikzpicture}%
Apellido y Nombre\quad\exPuntos{6.9cm}\par
\vspace{16pt}
Curso - Comisión - Especialidad: \exPuntos{4.6cm}\par

\vspace{44pt}
{\centering\fontsize{14.6}{17.5}\selectfont Examen Final Integrador\ \ \ - Emprendimientos Productivos\par}

\vspace{28pt}
\underline{Responda las siguientes preguntas:}\par

\vspace{28pt}
\exItem{1)}{¿Cuántas actividades puede ejercer simultáneamente un monotributista?}
\vspace{13pt}
\exItem{2)}{¿Cómo es el trámite de inscripción monotributista?}
\vspace{13pt}
\exItem{3)}{¿Es necesaria la clave fiscal?¿Y qué nivel requiere?}
\vspace{13pt}
\exItem{4)}{¿Cual es la diferencia entre venta de muebles y servicios?}
\vspace{13pt}
\exItem{5)}{¿Puede ser un monotributista de venta de servicios ser categoría K?}
\vspace{13pt}
\exItem{6)}{¿Qué es el convenio multilateral de IIBB? ¿Que es el monotributo Unificado?}
\vspace{13pt}
\exItem{7)}{¿Qué es el S.I.P.A?}
\vspace{13pt}
\exItem{8)}{¿Que cubre el aporte de la obra social del monotributo?}
\vspace{13pt}
\exItem{9)}{\ Ejercicio a desarrollar:}

\vspace{47pt}
\leftskip=36pt
Suponga que tiene un pequeño emprendimiento de \dotfill\par
\vspace{16pt}
y tiene una facturación mensualizada aproximada de \dotfill\par
\vspace{16pt}
tiene \exPuntos{2.9cm} locales físicos de venta al público, con las siguientes superficies\par
\vspace{30pt}
\exPuntos{4.2cm} y un consumo eléctrico mensual de \dotfill\ y\par
\vspace{16pt}
paga de alquiler \exPuntos{6.4cm} .\par
\leftskip=0pt

\vspace{57pt}
\exSub{a)}{¿Qué categoría de monotributo tendría que estar inscripto?}
\vspace{13pt}
\exSub{b)}{¿En qué actividades y qué códigos de actividades le corresponde según AFIP?}

% ---------------------------------------------------------- página 2
\newpage
\thispagestyle{empty}
\vspace*{50pt}
11) ¿Qué es el I.V.A?\par
\vspace{16pt}
12) ¿Cuáles son las distintas alícuotas de IVA?\par
\vspace{30pt}
12) Nombra algún producto que esté gravado con el 10,5\% de IVA.\par
\vspace{30pt}
13)¿Cuántos y cuáles tipos de facturación existen?\par
\vspace{30pt}
14) ¿Qué es la Factura A? ¿Quién la emite y hacia quién?\par
\vspace{30pt}
15) ¿Qué es la Factura B? ¿Quién la emite y hacia quién?\par
\vspace{30pt}
16) ¿Qué es la Factura C? ¿Quién la emite y hacia quién?\par
\vspace{30pt}
17) ¿Qué es la Factura E? ¿Quién la emite y hacia quién?\par
\vspace{30pt}
18) ¿Qué es la Factura T? ¿Quién la emite y hacia quién?\par
\vspace{30pt}
19) ¿Cuál es la ventaja impositiva de la Factura A?\par
\vspace{26pt}
20) ¿Qué significa las siglas C.U.I.T? ¿Qué implicancia tiene y cuál es su formato?\par

\endgroup
\clearpage
\restoregeometry
